\chapter{Peripheral Angioplasty and Stenting}

\section*{Overview}
Peripheral angioplasty and stenting is a minimally invasive procedure that opens narrowed or blocked arteries of the legs using a balloon and, when needed, a stent, to treat peripheral artery disease.

\section*{Alternative Names}
Lower-extremity angioplasty and stenting; peripheral vascular intervention; angioplasty of the leg arteries

\section*{Principle}
A balloon catheter is advanced over a guidewire through the arterial occlusion and inflated to dilate the vessel; a stent is placed to scaffold it open. This restores blood flow and relieves claudication or critical limb ischemia.

\section*{History}
Percutaneous transluminal angioplasty was developed by Andreas Gruentzig in the 1970s and applied to peripheral arteries soon after, with stents added in the 1990s.

\section*{Why It Is Done}
Performed for symptomatic peripheral artery disease (claudication, rest pain, or tissue loss) that is not controlled by medical therapy, and for acute limb ischemia from thrombosis or embolism.

\section*{Is This a Primary or Secondary Test or Procedure}
Primary or secondary revascularization for peripheral artery disease.

\section*{How It Is Performed}
Performed in an angiography suite under local anesthesia. Access is obtained (usually the femoral or popliteal artery), angiography maps the lesion, and a guidewire crosses it. The balloon is inflated and a stent is deployed as needed. Completion angiography confirms flow.

\section*{Preparation}
Preoperative imaging with duplex ultrasound, CT, or MR angiography; antiplatelet therapy; and assessment of renal function and contrast allergy.

\section*{Contraindications and Precautions}
Active infection, severe contrast allergy, and unreconstructable disease are relative contraindications. Precaution: calcified or long occlusions carry higher risk.

\section*{Risks}
Vessel injury or rupture, dissection, distal embolization, thrombosis, bleeding or pseudoaneurysm at the access site, and contrast nephropathy.

\section*{Aftercare and Recovery}
Observation, antiplatelet therapy, and access-site care. Most patients go home within a day.

\section*{Results and Interpretation}
Post-procedure imaging confirms a patent vessel with improved flow; clinical improvement in walking or rest pain confirms benefit.

\section*{Expected Results}
A patent, non-stenotic artery with restored distal flow and relief of symptoms.

\section*{Follow-up}
Duplex surveillance, risk-factor management, and repeat intervention for restenosis or progression.

\section*{References}
\begin{enumerate}
  \item MedlinePlus. Peripheral artery angioplasty and stenting. U.S. National Library of Medicine; 2025.
  \item Current Medical Diagnosis and Treatment. McGraw-Hill; 2025.
\end{enumerate}

\clearpage
